% font: footnotesize
% TAB-22: generated by harness/analysis/paper_outputs.py; do not edit (every cell's source is in TAB-22.cells.json)
\begin{tabularx}{\linewidth}{>{\hsize=1.218\hsize\raggedright\hyphenpenalty=10000\exhyphenpenalty=10000\arraybackslash}X>{\hsize=1.173\hsize\raggedright\hyphenpenalty=10000\exhyphenpenalty=10000\arraybackslash}X>{\hsize=0.609\hsize\raggedright\hyphenpenalty=10000\exhyphenpenalty=10000\arraybackslash}X}
\toprule
Stamped field & Stamped value (or \mbox{[a description]}) & In the paper \\
\midrule
topup: costs.price\_table\_date & \mbox{2026-10-12} & The analysis code's own price-table date; the run's price date is \mbox{2026-10-13}, as our ruling of \mbox{2026-10-13} records. The prices it applies are the same list prices the run was cleared on, re-checked unchanged before the run; the paper states list prices once, with their check date. \\
\addlinespace
topup: inputs.plan\_sha256 & \mbox{[a SHA-256 hash, in full in Panel~1]} & The driver was given the run folder's own copy of the plan (topup\_plan.json), so its plan check compares that copy with itself; the value is that copy's file hash, as in the run's row of our cost records. Before the driver ran, the cleared plan's hash (e038a46e\ldots{}, as in that row) and all 13 run-file hashes (the top-up's and the main run's, as in their rows of our cost records) were checked (our work log's entry for the driver run). \\
\addlinespace[1em]
main: ``reporting.group\_b.item\_9b.arms.rt-haiku-med:1x6.batch\_mean'' & \mbox{[equals its list mean; its batch\_label reads `real time (list)']} & The six real-time settings were billed at list price, so their `batch\_mean' holds the list cost; the paper reports them at list. \\
\bottomrule
\end{tabularx}
